% Generated draft: requires mathematical and source review.
\documentclass[11pt]{article}
\usepackage[margin=25mm]{geometry}
\usepackage{fontspec}
\setmainfont{Latin Modern Roman}
\usepackage{amsmath,amssymb,mathtools,unicode-math}
\setmathfont{Latin Modern Math}
\usepackage{enumitem,needspace,xurl,hyperref}
\setlength{\parindent}{0pt}
\setlength{\parskip}{5pt}
\setlength{\emergencystretch}{3em}
\newcommand{\N}{\mathbb{N}}
\newcommand{\Z}{\mathbb{Z}}
\newcommand{\Q}{\mathbb{Q}}
\newcommand{\R}{\mathbb{R}}
\newcommand{\C}{\mathbb{C}}
\newcommand{\F}{\mathbb{F}}
\newcommand{\A}{\mathbb{A}}
\newcommand{\PP}{\mathbb{P}}
\newcommand{\E}{\mathbb{E}}
\newcommand{\eps}{\varepsilon}
\newcommand{\dd}{\,\mathrm{d}}
\begin{document}
\section*{UnsolvedMath ID: 700050}
\textbf{Source problem code:} catalogue entry
\subsection{Definitions and mathematical statement}
Let \(\kappa\) denote the prescribed number of free generators, and let \(F_\kappa\) be a free group with a free basis of cardinality \(\kappa\). The supplied record does not specify the permitted values of \(\kappa\).

A strict total order \(<\) on \(F_\kappa\) is a transitive, irreflexive relation such that, for distinct \(g,h\in F_\kappa\), exactly one of \(g<h\) and \(h<g\) holds. Under the two-sided compatibility convention, a linearly ordered group satisfies
\[
\forall a,b,g,h\in F_\kappa,\qquad g<h\ \Longrightarrow\ agb<ahb.
\]
Under the left-order convention, compatibility instead means
\[
\forall a,g,h\in F_\kappa,\qquad g<h\ \Longrightarrow\ ag<ah.
\]
The supplied record does not explicitly specify the compatibility convention; these alternatives must not be identified without clarification.

\textbf{Classification problem (A. I. Kokorin, 1965).} Give an elementary classification of linearly ordered free groups with a fixed number \(\kappa\) of free generators.

One possible classification target is the collection of compatible strict total orders on \(F_\kappa\), up to order-preserving group isomorphism. For this target, two orders \(<_1\) and \(<_2\) are equivalent precisely when there exists \(\varphi\in\operatorname{Aut}(F_\kappa)\) such that
\[
\forall g,h\in F_\kappa,\qquad g<_1h\ \Longleftrightarrow\ \varphi(g)<_2\varphi(h).
\]
However, the source does not specify the classification equivalence, whether a chosen free basis must be preserved, or what constitutes an elementary classification. This formulation therefore records a possible interpretation rather than resolving those qualifications.
\subsection{Short English statement}
For a fixed number of free generators, classify the linear group orders on the corresponding free group by elementary means. The source leaves the order-compatibility convention, permitted ranks, classification equivalence, and meaning of ``elementary'' unspecified. Its open status is a supplied registry label, not an independently verified assessment.
\subsection{Sources}
\begin{itemize}
\item Local UnsolvedMath registry: supplied complete source record, permanent ID \(700050\), exact source problem code \texttt{catalogue entry}, title ``Kourovka Notebook Problem 1.55.'' The record labels the problem \texttt{open}; its creation and update timestamps are both \texttt{2026-10-16T14:00:00Z}. These are supplied registry data, not independent mathematical or literature verification.
\item A. I. Kokorin, Problem 1.55, \emph{The Kourovka Notebook (Unsolved Problems in Group Theory)}, first issue (1965), as attributed in the supplied record. The record identifies the twenty-first issue edition (2026), arXiv:1401.0300v48: \url{https://arxiv.org/abs/1401.0300v48}.
\item The supplied collection description attributes editorship of \emph{Kourovka Notebook (Issues 1--21)} to E. I. Khukhro and V. D. Mazurov. Supplied upstream revision: \texttt{d371ee917ac9647fbb06cd6fc529deec6f3dda82}. The cited edition has not been independently consulted for this draft.
\end{itemize}
\end{document}
